\section{Introduction}
\IEEEPARstart{L}{arge} language-model agents are being explored for
plant-facing tasks. Xia et al. present an end-to-end industrial-control
framework using structured prompts and event-driven production context
\cite{xia2025control}. Other work studies operator-reviewed control-sequence
generation \cite{huang2026made}, modular factory skills
\cite{koehle2026llmdriven}, and production scheduling
\cite{waseem2026pretrained}. Together, these studies establish a broad
industrial-agent research direction, with different task and integration
boundaries. Across plants, agent roles, signals, protocols, objectives, and
production context vary. This paper focuses on the software integration
problem of configuring those differences around reusable task services,
registered harnesses, shared operation checks, semantic process I/O, and
route-specific records; it does not claim the first agent-to-control
integration.

Consider injection molding. An agent team may use recent measurements and
process knowledge to diagnose insufficient packing and propose a holding-
pressure change. To carry that proposal into an auditable production task, a
system must resolve the active line, product, recipe, and run; map the
quantity to the correct device interface; check that the requested operation
is available for that context; and retain the command evidence alongside
subsequent quality observations. Another process may use different signals,
protocols, or objectives while reusing the same collaboration and runtime
services. The design problem is therefore to make scene-specific variation
configurable while keeping each operation tied to its production context.

This paper presents AgentWorkShop, a configurable industrial-agent framework
that connects collaborative tasks and registered harnesses to production
context, manufacturing execution system (MES) data, analysis and modeling
services, data-acquisition (DAQ) and write-control (DCW) nodes, and
protocol-specific drivers. Its access surfaces include a command-line
interface (CLI), web, software development kit (SDK), representational state
transfer (REST), Model Context Protocol (MCP), and Agent-to-Agent (A2A)
interfaces.
Scene blueprints define signal maps, objectives, and guards; channel profiles,
tool policies, and agent bindings scope operations. Agents and callers invoke
named operations through shared runtime checks; the framework does not choose
routes autonomously and retains proposals, command confirmation, and process
observations as distinct records. Figure~\ref{fig:arch} presents the system
architecture, and Fig.~\ref{fig:runtime} shows the cross-module task flow.
Section~IV details the node-exploration and recipe lifecycles.

The paper makes three contributions:
\begin{enumerate}
\item It formulates five deployment requirements and six design rules for
  connecting agent collaboration, production context, process interfaces, and
  route-specific evidence. These are explicit design choices for the target
  workflows, not a taxonomy claimed to be empirically derived.
\item It specifies AgentWorkShop's cross-module execution flow and software
  modules, including task and harness integration, production-context
  configuration, analysis and modeling services, semantic device access, and
  route-specific evidence handling.
\item It reports complementary evidence from historical simulator and
  protocol archives, live-engine tasks, a deterministic controller,
  deployment regressions, and a recipe-path case. A current-version
  application-level two-scene integration run verifies protocol reads, bounded writes,
  readback, and data lineage, while its model gates reject both candidate
  datasets. Each evidence unit is tied to its interface and software
  version; the results do not establish successful current-version
  optimization, multi-harness portability, or comparative control benefit.
\end{enumerate}

Section~II reviews related industrial-agent and integration work.
Section~III presents the design rationale, architecture, and cross-module
execution flow; Section~IV explains module responsibilities and their
integrated operation. Section~V describes the experiments and results, and
Sections~VI and~VII discuss the evidence limits and conclude the paper.
